%=====================================================================
% 附录  Hitchin 模空间
% 原文：Appendix. Hitchin's moduli spaces（arXiv:1702.06499v2, 附录）
%=====================================================================
% 原文宏 \BC（即复数域 \mathbb{C}）：若共享 preamble 已提供则此处自动跳过，
% 未提供时就地补定义，以保证公式逐字无损且整篇可编译。
\providecommand{\BC}{\mathbb{C}}

\section{Hitchin 模空间}\label{Gloss}

\begin{keytakeaways}
\begin{itemize}
  \item \textbf{物理图像：} Hitchin 模空间 $\CM_H(G)$ 是全文三条进路共同的目标流形兼相空间。Kapustin--Witten 的进路把它取作 $N=4$ 超 Yang--Mills 理论沿 $C$ 约化后所得二维 sigma 模型的目标流形；AGT 对应及其 Nekrasov--Shatashvili（NS）极限把它看作被量子化的 Hitchin 可积系统的相空间，共形块给出 Hitchin 哈密顿量的本征函数；Beilinson--Drinfeld 的 CFT 进路则经由 $\CM_{dR}$ 中的 oper 子空间与它对接。
  \item \textbf{数学工具：} 谱曲线 $\Sigma\subset T^*C$ 连同 Hitchin 纤维化把 $\CM_H(G)$ 实现为底空间 $\CB\simeq H^0(C,K^2)$ 上的代数完全可积系统；非阿贝尔 Hodge（NAH）对应给出 $\CM_H(G)$、$\CM_{dR}(G)$ 与 $\CM_B(G)$ 三种描述的同一性；以 $\mathbb{P}^1$ 参数化的一族复结构 $I_\ze$ 与全纯辛形式 $\Om_\ze$ 构成超 K\"ahler 结构，是统一上述描述的语言。
  \item \textbf{关键结论：} 同一流形上的复结构 $I$、$J$、$K$ 分别对应 Higgs 丛侧、平坦联络（de Rham）侧与基本群表示（Betti，即特征簇）侧。三条进路的差别恰好体现为取用了族中不同的复结构——Kapustin--Witten 的 Hecke 本征膜按复结构 $I$ 下的 Hitchin 纤维化描述，Beilinson--Drinfeld 一侧取的是 $\CM_{dR}$ 中的局部系统——而由 $\vf\ra e^{i\theta}\vf$ 生成的超 K\"ahler 旋转（在 $\ze\neq0,\infty$ 处可由 $\CM_H(G)$ 的微分同胚实现）把三者联系起来。
  \item \textbf{附录的作用：} 正因如此，Hitchin 模空间才是三条进路可以逐项比较的公共几何框架；本附录以术语表形式固定全文记号，条目之间的衔接（谱曲线—纤维化、NAH—三种描述、$\Om_\ze$—辛结构）正是正文论证的几何骨架。
\end{itemize}
\end{keytakeaways}

\bigskip
\noindent\rule{\textwidth}{0.6pt}
\bigskip

本文假定 $G=SL(2)$，且 $C$ 是亏格为 $g$、带 $n$ 个穿刺点的 Riemann 面。

{\bf Hitchin 模空间 $\CM_H(G)$} \cite{Hi87}。
这是偶 $(\CE,\vf)$ 的模空间，其中 $\CE=(E,\bar{\partial}_{\CE})$ 是光滑向量丛 $E$ 上的全纯结构（holomorphic structure），而 $\vf\in H^0(C,\mathrm{End}(\CE)\ot K)$。这类偶模去自然规范变换后所得的模空间记作 $\CM_H(G)$。

{\bf Hitchin 可积系统} \cite{Hi87}。
给定 $(\mathcal{E},\vf)$，可构造谱曲线
$\Sigma=\{(u,v);v^2=\frac{1}{2}\mathrm{tr}(\vf^2)\}\subset T^*C$，
以及代表 $\vf-v$ 的余核的线丛 $\CL$。
由 $(\Sigma, \mathcal{L})$ 可以重构 $(\mathcal{E},\vf)$：
$\mathcal{E}=\pi_*(\mathcal{L})$，$\vf=\pi_*(v)$。
这就把 $\CM_{H}(G,C)$ 描述为底空间
$\CB\simeq H^0(C,K^2)$
上的环面纤维化，其中纤维——即 $\CL$ 的选取——当 $G=GL(2)$ 时等同于 $\Sigma$ 的 Jacobi 簇（Jacobian），当 $G=SL(2)$ 时等同于 Prym 簇（Prym variety）。底空间 $\CB$ 的自然坐标由 Hitchin 哈密顿量给出，其定义为将
$\frac{1}{2}\mathrm{tr}(\vf^2)=\sum_{r=1}^{3g-3+n}\vartheta_rH_r$
展开，其中 $\{\vartheta_r, r=1,\dots,3g-3+n\}$ 构成 $H^0(C,K^2)$ 的一组基。

{\bf 局部系统}。即偶 $(\CE,\nabla_\ep')$，其中 $\CE$ 是如上全纯向量丛，$\nabla_\ep'$ 是全纯 $\ep$-联络，对函数 $f$ 与 $E$ 的光滑截面 $s$ 满足
$\nabla_\ep' (fs)=\ep (\pa f)s+f\nabla_\ep's$。这类偶的模空间记作 $\CM_{dR}(G)$。
此处常把局部系统等同于相应的\emph{平坦丛}（flat bundle），即转移函数为常值的局部平凡化系统；或等同于由 $(\CF,\nabla_\ep')$ 的和乐得到的基本群表示（模去共轭）。由此得到 $\CM_{dR}(G)$ 与

{\bf 特征簇（character variety）$\CM_B(G)$}：$\pi_1(C)$ 到 $G$ 的表示模去整体共轭所成之空间；作为代数簇，它被描述为一个 GIT 商
$\BC[\mathrm{Hom}(\pi_1(C),G]^G$。

{\bf oper}。这是一类特殊的局部系统：其中 $\CE=\CE_{\rm op}$，即唯一可使全纯联络 $\nabla_\ep'$ 取形式
$\nabla_\ep'=dz\big(\ep\pa_z+\big(
\begin{smallmatrix} 0 & u \\ 1 & 0\end{smallmatrix}\big)\big)$
的扩张
$0\ra K^{1/2}\ra\CE_{\rm op} \ra K^{-1/2}\ra 0$。

{\bf 非阿贝尔 Hodge（NAH）对应} \cite{Hi87,Si}。给定 Higgs 偶 $(\CE,\vf)$，存在
唯一的调和度量 $h$ 满足
%\begin{align}\label{selfdual}
$F_{\CE,h}+R^2[\vf,\vf^{\dagger_h}]=0$
%\end{align}
其中 $F_{\CE,h}$ 是唯一以
$(0,1)$-部分为 $\bar{\partial}_{\CE}$ 的 $h$-酉联络 $D_{\CE,h}$ 的曲率。
% In the holomorphic gauge where $\bar{\partial}_{\CE}=\bar{\partial}$,
%we have $D_{\CE,h}=\partial+ h(\partial h^{-1}) +\bar\partial$, and \rf{selfdual} becomes
%\begin{align}\label{selfdual'}
%\bar\partial(h(\partial h^{-1}))+R^2\big[\vf,\bar{h}^{-1}\vf^{\dagger}\bar{h}\big]=0\,,
%\end{align}
%Equation \rf{selfdual} determines $h$ uniquely for given $(\CE,\vf)$.  
于是可以构造相应的双参数平坦联络族
%\begin{align}
$\nabla_{\zeta,R}=\zeta^{-1}{R}\,\vf+D_{\CE,h}+R\zeta\,\vf^{\dagger_h}$。
%\end{align}
把 $\nabla_{\zeta,R}$ 分解为 $(1,0)$ 与 $(0,1)$-两部分 %$\nabla'$ and $\nabla''$ 
即给出一个偶 $(\CF,\nabla_\ep')$，它由
$\CF=(E,\bar{\partial}_{\CF})$ 与 $\ep$-联络
$\nabla_\ep'=\ep\nabla'=\ep\pa_{\CE,h}+\vf$ 组成，其中 $\ep=\zeta/R$，且 $\nabla_\ep'$ 在
$\bar{\partial}_{\CF}$ 所定义的复结构下全纯。

{\bf 超 K\"ahler 结构} \cite{Hi87}。存在由 $\mathbb{P}^1$ 参数化的一族复结构 $I_\zeta$ 与全纯
辛结构 $\Omega_\zeta$。后者定义为
$
\Omega_\zeta =\frac{1}{2}\int_C\;{\rm tr}(\de\mathcal{A}_\zeta\wedge\de\mathcal{A}_\zeta)
$.
三重辛形式 $(\omega_I,\omega_J,\omega_K)$ 可由 $\Omega_\zeta$ 的展开定义：
$
\Omega_\zeta=\frac{1}{2\zeta}(\omega_J+i\omega_K)+i\omega_I+\frac{1}{2}{\zeta}(\omega_J-i\omega_K).
$
相应的复结构为
$
I_{\zeta}=\frac{1}{1+|\zeta|^2}((1-|\zeta|^2)I-i(\zeta-\bar{\zeta})J-(\zeta+\bar{\zeta})K).
$

{\bf 复 Fenchel--Nielsen 坐标} \cite{NRS}。这是与 $C$ 的裤子分解（pants decomposition）$\si$ 相对应的 $\CM_{B}(G)$ 上的 Darboux 坐标，而该分解由沿闭曲线 $\ga_i$（$i=1,\dots,3g-3+n$）切割而得。
复长度坐标把迹函数 $L_i=\mathrm{tr}(\rho(\ga_i))$ 参数化为 $L_r=2\cosh(a_r/2)$。可以定义正则共轭坐标
$\kappa_r$，使得自然的 Poisson 结构表示为
$\{a_r,\kappa_s\}=\de_{r,s}$、$\{a_r,a_s\}=0=\{\kappa_r,\kappa_s\}$。

\bigskip
\noindent\rule{\textwidth}{0.6pt}
% [Tutor Note: 被注释掉的三个公式环境（原文写作 %\begin{align}…%\end{align}）内容如下：其一是 Hitchin 1987 的核心方程 $F_{\CE,h}+R^2[\vf,\vf^{\dagger_h}]=0$，即 Hitchin 方程，亦称自对偶方程（它是 4d 自对偶 Yang–Mills 方程在 $C$ 上的 $G$-不变约化，故原文径称 self-dual）；其二给出全纯规范 $\bar\partial_{\CE}=\bar\partial$ 下解 $h$ 的等价非线性 PDE $\bar\partial(h(\partial h^{-1}))+R^2[\vf,\bar h^{-1}\vf^\dagger\bar h]=0$；其三是 $\nabla_{\zeta,R}$ 的 align 环境。它们被注释的原因应是排版遗留而非数学内容删改：同一段里被引用的 \label{selfdual}、\label{selfdual'} 也一并被注释，注释中的 \rf{selfdual}（=\eqref）因此指向不存在的 label，作者最终改以行内无编号公式呈现。译者严格保持注释状态、不改回可见公式，并按“公式逐字保留”原则保留其行内可见版本。]
% [Tutor Note: $R^2$ 项的来源：$R$ 是 4d$\to$2d 约化中紧化曲面 $C$ 的尺度／耦合参数（Hitchin 原文即以 $R$ 记之）；本文正文 Conformal limit 一节明确它由 NAH 对应中的标度 $\vf\ra R\vf$ 引入，并把保持 $\ze/R$ 有限的极限 $R\ra0$、$\ze\ra0$ 作为几何 Langlands 的关键。核对量纲：$F_{\CE,h}$ 与 $[\vf,\vf^{\dagger_h}]$ 都是 $(1,1)$-形式（$\vf$ 为伴随值 $(1,0)$-形式、取值于 $K$，$\vf^{\dagger_h}$ 为 $(0,1)$-形式），故 $R^2$ 是使两项同阶的无量纲因子；按 $\vf\ra R\vf$ 改写时第二项整体乘 $R^2$，与正文标度约定一致。注意 $\vf^{\dagger_h}$ 显含 Hermitian 度量 $h$，这正是该方程作为 $h$ 的方程、而非代数方程的原因。]
% [Tutor Note: 复结构 $I,J,K$ 与 $\Om_\ze$ 的约定核对：由附录的 $\Om_\ze$ 展开可得 $\ze=0$ 时 $I_\ze=I$、$\ze=i$ 时 $I_\ze=J$、$\ze=-1$ 时 $I_\ze=K$（代入 $1+|\ze|^2$：$\ze=i$ 给出 $I_i=\frac12(-i\cdot2i)J=J$；$\ze=-1$ 给出 $I_{-1}=\frac12\cdot2K=K$）。注意 $\Om_\ze$ 在 $\ze=0$ 有极点，故“$\ze=0$ 处的全纯辛形式”应理解为 $\ze\Om_\ze|_{\ze=0}=\frac12(\omega_J+i\omega_K)$，这才是 $I$-全纯的 $(2,0)$-形式；文献对 $\ze$ 是否吸收这个 $1/\ze$ 归一化不同，比较公式前须核对。另：$\de$ 是模空间（场空间）上的外微分，$\mathcal{A}_\ze$ 为 $\CM_H$ 上对应的联络形式，$\wedge$ 同时含 $C$ 上的形式与场空间微分，$\Om_\ze$ 才是 $\CM_H$ 上的辛形式。]
% [Tutor Note: 稳定性条件与“模空间”的确切含义：原文只说“模去自然规范变换”（modulo natural gauge transformations），但 $\CM_H(G)$ 的严格定义须取 $(E,\vf)$ 的 S-等价类，即 polystable 对象。Hitchin 方程可解 $\Leftrightarrow$ $(E,\vf)$ polystable，此即 Hitchin (1987) 与 Simpson (1992) 的定理；附录中 NAH 对应断言的 $\CM_H(G)\cong\CM_{dR}(G)\cong\CM_B(G)$ 实为 polystable 模空间之间的同构，稳定性条件不可省略。另注意 $\vf$ 是伴随值 $(1,0)$-形式（见 $\mathrm{End}(\CE)\ot K$），$\mathrm{tr}(\vf^2)\in H^0(C,K^2)$ 是二次微分而非函数；对 $G=SL(2)$ 还应取无迹部分 $\mathrm{End}_0$，否则底空间会多出 $H^0(C,K)$ 方向（详见交付报告“原文可疑之处”第 2 条）。]
% [Tutor Note: 谱曲线与 Hitchin 纤维化：$\Sigma=\{(u,v);v^2=\frac12\mathrm{tr}(\vf^2)\}\subset T^*C$ 是 $\vf$ 的特征方程定义的二重覆盖（$v$ 为 $T^*C$ 上的 tautological 1-形式），$\CL=\mathrm{coker}(\vf-v)$ 是 $\Sigma$ 上的线丛；$\CE=\pi_*\CL$、$\vf=\pi_*(v)$ 即“Higgs 丛 = 谱曲线上线丛的直像”。$G=SL(2)$ 时 $\det\CE$ 平凡，故 $\CL$ 须满足 $\mathrm{Nm}(\CL)=\CO_C$，其纤维是 Prym 簇，即 $\ker(\mathrm{Nm}:\mathrm{Pic}(\Sigma)\to\mathrm{Pic}(C))$ 的单位连通分支，维数 $g_\Sigma-g=3g-3+n$，恰等于底空间 $\CB$ 的维数，故 $\dim_\BC\CM_H=2(3g-3+n)$。Hitchin 哈密顿量 $H_r$ 是 $\frac12\mathrm{tr}(\vf^2)$ 在基 $\vartheta_r$ 下的系数；正文 NS 极限中 oper 的 $u=u^\sharp+\sum_rE_r\vartheta_r$ 用的正是同一组 $\vartheta_r$，$E_r$ 为 $D_r\Psi=E_r\Psi$ 的本征值——附录与正文的 $\vartheta_r$ 记号须一并理解。]
